
\newcommand{\Title}{Proposal Title (see preamble.tex)}
\usepackage[breaklinks=true, pdftex, 
  pdfborder={0 0 0}, colorlinks=true, linkcolor=blue, citecolor=blue, urlcolor=blue, 
  bookmarks=false, pdfpagelabels=false, pdftitle={\Title}]{hyperref}
\usepackage{fullpage}
\usepackage{times}
\usepackage{ifthen}   % for conditionals
\usepackage{graphicx,graphics}
\usepackage{sidecap}
\usepackage[]{rotating}
\usepackage{setspace}
\usepackage{amssymb,amsmath,amsthm}
\usepackage{color}
\usepackage{multirow}
\usepackage{pdfpages}
\usepackage{longtable} 	%used in Masha's c&p
\usepackage{datetime}
\usepackage{xspace}
\usepackage{enumitem} % For tighter lists
\usepackage{refcheck}           % For checking references
\usepackage[it,small]{caption}
\usepackage[top=1in,bottom=1in,left=1in,right=1in]{geometry} % Set margins
\usepackage{subcaption,wrapfig}
\usepackage[utf8]{inputenc}
\usepackage[T1]{fontenc}
\usepackage{doi}
\usepackage[numbers,sort&compress]{natbib}
\usepackage{algorithmic}
% \usepackage[noresetcount,lined,boxed]{algorithm2e} % ... for algorithms
\usepackage[linesnumbered,lined,boxed]{algorithm2e} % ... for algorithms
\newlength{\commentWidth}
\setlength{\commentWidth}{2in}
\newcommand{\atcp}[1]{\tcp*[r]{\makebox[\commentWidth]{#1\hfill}}}

\graphicspath{{./figs/}} 
\usepackage{pgfgantt}
\usepackage{pdflscape} % provides the landscape environment
%\usepackage{mathtools}
\usepackage[most]{tcolorbox}
\usepackage{textcomp}

\usepackage{titlesec}
%% %%% The following is a fix for people using the buggy version of titlesec
  \usepackage{etoolbox}
  \makeatletter
  \patchcmd{\ttlh@hang}{\parindent\z@}{\parindent\z@\leavevmode}{}{}
  \patchcmd{\ttlh@hang}{\noindent}{}{}{}
  \makeatother
% %%%
\titleformat{\subsubsection}[runin]{\bfseries}{\thesubsubsection}{1em}{}[.~~]
  \titlespacing{\section}{0pt}{*1.35}{*0.35}
  \titlespacing{\subsection}{0pt}{*1.15}{*0.15}
  \titlespacing{\subsubsection}{0pt}{*1.05}{*0.05}
 \titlespacing{\paragraph}{0pt}{*1.15}{*0.5}

% define highlight box style
\usepackage{mdframed}                 % creating color text boxes
\usepackage{colortbl}
% \usepackage[table,dvipsnames]{xcolor}  
\usepackage[most]{tcolorbox}
\usepackage{enumitem}

%\usepackage{xcolor}
%\definecolor{highlight}{HTML}{cdeeff} % background color
%%\definecolor{highlight}{HTML}{aed6f1} % background color
\definecolor{optim}{HTML}{DFEFF9}
\mdfdefinestyle{highlightoptim}{leftmargin=0pt,rightmargin=0pt,%
innerleftmargin=3pt,innerrightmargin=2pt,%
innertopmargin=4pt,innerbottommargin=4pt,%
linewidth=2pt,%
backgroundcolor=blue!10!white,linecolor=blue!50!black}
\definecolor{stats}{HTML}{D4EFDF} % background color
\mdfdefinestyle{highlightstats}{leftmargin=0pt,rightmargin=0pt,%
innerleftmargin=3pt,innerrightmargin=2pt,%
innertopmargin=4pt,innerbottommargin=4pt,%
linewidth=2pt,%
backgroundcolor=green!10!white,linecolor=green!50!black}
\definecolor{appl}{HTML}{F2D7D5} % background color
\mdfdefinestyle{highlightsample}{leftmargin=0pt,rightmargin=0pt,%
innerleftmargin=3pt,innerrightmargin=2pt,%
innertopmargin=4pt,innerbottommargin=4pt,%
linewidth=2pt,%
backgroundcolor=red!10!white,linecolor=red!50!black}
\definecolor{AD}{HTML}{FFFFCC} % background color
\mdfdefinestyle{highlightsoft}{leftmargin=0pt,rightmargin=0pt,%
innerleftmargin=3pt,innerrightmargin=2pt,%
innertopmargin=4pt,innerbottommargin=4pt,%
linewidth=2pt,%
backgroundcolor=yellow!10!white,linecolor=yellow!50!black}

%%Space things:
\setlength{\textfloatsep}{2pt plus 1.0pt minus 1.0pt}
\setlength{\floatsep}{2pt plus 1.0pt minus 1.0pt}
\setlength{\intextsep}{2pt plus 1.0pt minus 1.0pt}
\setlength{\belowcaptionskip}{1pt}
\setlength{\abovecaptionskip}{1pt}

% Various codes (for consistency and to differentiate from acronyms)
\newcommand{\R}{\mathbb R}           % Real  numbers
\newcommand{\Z}{\mathbb Z}           % Integer  numbers
\newcommand{\E}{\mathbb E}           % expectation
\newcommand{\Prob}{\mathbb P} %probability
\newcommand{\Risk}{\mathbb R}           % expectation
\newcommand{\dps}{\displaystyle}
\newcommand{\maxi}{\mathop{\mbox{maximize}}}
\newcommand{\mini}{\mathop{\mbox{minimize}}}
\newcommand{\amin}{\mathop{\mbox{argmin}}}
\def\IhH {I_h^H}
\def\IhHb { {\hat I}_h^H}
\newcommand{\DOESC}{DOE/SC} % US Dept of Energy (DOE) Office of Science (SC)
\newcommand{\DoE}{DoE} % Design of Experiments (could have many modifiers)
\newcommand{\st}{\mbox{subject to}}
\newcommand{\stn}{\mbox{s.t.}}
\newcommand{\bvec}{\left(\begin{array}{c}}
\newcommand{\evec}{\end{array}\right)}
\newcommand{\cF}{{\cal F}}
\newcommand{\cS}{{\cal S}}
\newcommand{\cP}{{\cal P}}
\newcommand{\xh}{\hat{x}}
\newcommand{\yh}{\hat{y}}
\newcommand{\Lag}{{\cal L}}
\newcommand{\sLag}{{\tilde{\cal L}}}
\newcommand{\sI}{{\tilde{I}}}
\newcommand{\heta}{\hat{\eta}}
%\newcommand{\homega}{\hat{\omega}}
\newcommand{\hkappa}{\hat{\kappa}}
\newcommand{\x}{\mathbf x}
\newcommand{\y}{\mathbf y}
\newcommand{\cc}{\mathbf c}

\DeclareGraphicsRule{.pdftex}{pdf}{.pdftex}{}

% ... theorems etc
\newtheorem{definition}{Definition}[section]
\newtheorem{theorem}{Theorem}[section]
\newtheorem{lemma}{Lemma}[section]
\newtheorem{proposition}{Proposition}[section]
\newtheorem{assumption}{Assumption}[section]
\newtheorem{corollary}{Corollary}[section]
\newtheorem{property}{Property}[section]
\newtheorem{remark}{Remark}[section]

% Commands specific to this proposal
\newcommand{\subsubsubsection}[1]{\paragraph{#1.}}
%{\vspace{1pt} \noindent \textbf{#1}}
%\newcommand{\mcomment}[2]{\textcolor{red}{#2}\marginpar{{\color{red}\textbf{#1 Note}}}}
\newcommand{\call}[1]{\noindent\textcolor{blue}{\textbf{From the call:} #1}}
%\renewcommand{\call}[1]{ } % Final version
 \newcommand{\team}[1]			% team command
      {\textcolor{cyan}{#1}\marginpar{{\color{cyan}\textbf{Team}}}}
% \newcommand{\team}[1]{}			% team command
  
\newcommand{\ignore}[1]{}  % Use to comment out large blocks of text
\newcommand{\BigO}[1]{\ensuremath{\mathcal{O}\!\left(#1\right)}}
\newcommand{\Sec}[1]{{\S\ref{#1}}}
%\newcommand{\Secs}{Secs.}
\newcommand{\Fig}[1]{Fig.~\ref{#1}}
%\newcommand{\Figs}{Figs.}


\newcommand{\pagebudget}[1]{\textcolor{red}{~#1}} % For working version
%\renewcommand{\pagebudget}[1]{} % In final version


% For using headers ********************************
\usepackage{fancyhdr}
\pagestyle{fancy}
\renewcommand{\headrulewidth}{0pt}
\renewcommand{\footrulewidth}{0pt}
\fancyhead{}
\fancyhead[R]{{\it \rightmark}}
\fancyhead[L]{{\it \leftmark}}
\renewcommand{\sectionmark}[1]{\markright{#1}{}}
\voffset= -65pt
\headsep= 51pt %\voffset-12pt
\headheight=14pt
% headers *******************************************

%  Following not defined in  LaTeX2e <2003/12/01> - Steve
% \pdfminorversion=6 % Turns of PDF version warnings


% \newcommand{\bio}[3]{\newpage \phantomsection 
% 		 \pagestyle{fancy} \sectionmark{BIOGRAPHICAL SKETCHES \hfill #1}
%    \IfFileExists{bio/#3.tex}{\addcontentsline{toc}{subsection}{#1 -- #2} \input{bio/#3}
%     }{\IfFileExists{bio/#3.pdf}{\addcontentsline{toc}{subsection}{#1 -- #2}
%          \includepdf[pages=-,offset=0 -65,pagecommand={\thispagestyle{fancy}}]{bio/#3}
%       }{% If the desired PDF file doesn't exist, warn about it but continue
% 	 \addcontentsline{toc}{subsection}{\textcolor{red}{#1 -- #2 MISSING}}
%          \begin{center}
%             \textbf{\color{red}\Large No file found for `#1'!\\
%                     Tried bio/#3.tex and bio/#3.pdf} \\
%          \end{center}}}}

% \newcommand{\cnp}[3]{\newpage \phantomsection 
% 		 \pagestyle{fancy} \sectionmark{SUPPORT OF INVESTIGATORS \hfill #1}
%    \IfFileExists{cnp/#3support.tex}{\addcontentsline{toc}{subsection}{#1 -- 
% #2} \input{cnp/#3support.tex}
%     }{\IfFileExists{cnp/#3.pdf}{ 
% \addcontentsline{toc}{subsection}{#1 -- #2}
%          \includepdf[pages=-,offset=0 -65,pagecommand={\thispagestyle{fancy}} 
% ]{cnp/#3.pdf}
%       }{% If the desired PDF file doesn't exist, warn about it but continue
% 	 \addcontentsline{toc}{subsection}{\textcolor{red}{#1 -- #2 MISSING!!!}}
%          \begin{center}
%             \textbf{\color{red}\Large No file found for `#1'!\\
%                     Tried cnp/#3.tex and cnp/#3.pdf} \\
%          \end{center}}}}

% \newcommand{\todd}[1]{{{\color{blue}{Todd:} #1}}}
% \newcommand{\sven}[1]{{{\color{red}{Sven:} #1}}}

\usepackage[normalem]{ulem} % For \sout
\newcommand\hcancel[2][black]{\setbox0=\hbox{$#2$}\rlap{\raisebox{.45\ht0}{\textcolor{#1} {\rule{\wd0}{1pt}}}}#2} 
\newcommand{\replace}[2]{{\color{red}\sout{#1}\color{black}{\color{red}#2\color{black}}}} %TeX source markup.
\newcommand{\replacemath}[2]{{\hcancel[red]{#1}{}{\color{red}#2\color{black}}}} %TeX source markup.
\newcommand{\slnote}[1]{{{\color{blue}{ SL note:} #1}}\marginpar{\color{blue}{SL}}}
\newcommand{\mcnote}[1]{{{\color{cyan}{ MC note:} #1}}\marginpar{\color{cyan}{SP}}}
\newcommand{\tlnote}[1]{{{\color{brown}{ TL note:} #1}}\marginpar{\color{brown}{MM}}}
%% \renewcommand{\slnote}[1]{}
%% \renewcommand{\mcnote}[1]{}

\newcommand{\SLreverse}[2]{{\color{purple}\sout{#2}\color{black}{\color{purple}#1\color{black}}}} %TeX source REVERSEmarkup.
%\renewcommand{\SLreverse}[2]{#1}{}

\newcommand{\SLreplace}[2]{{\color{blue}\sout{#1}\color{black}{\color{blue}#2\color{black}}}} %TeX source markup.
%\renewcommand{\SLreplace}[2]{#2}{}

\usepackage{adjustbox}

\newcommand{\prtl}[2]{\frac{\partial #1}{\partial #2}}

\usepackage{xpatch}
\xpretocmd{\eqref}{Eq.~}{}{}

\usepackage{textpos}
\NewEnviron{RandSample}{%
	\colorbox{green!15!white}{
		\parbox{0.97\hsize}{
			\begin{textblock*}{75pt}[1, 1](\hsize,5pt)
			\emph{$\gets$ Sampler}
			\end{textblock*}%
			\BODY
		}
	}
}

\usepackage{tocloft}
\setlength{\cftsecnumwidth}{2em} 
